\chapter{Análise léxica baseada em expressões regulares}

\section{Introdução}

No capítulo anterior, apresentamos um analisador léxico ad-hoc para
uma linguagem simples. Ao tentarmos estender a linguagem considerada
com construções simples com identificadores ou comentários,
encontraremos problemas: temos que modificar a estrutura do estado do
analisador léxico para acomodar uma string para acumular caracteres
que formam um possível identificador e modificar toda a estrutura do
analisador para lidar com comentários de linha / blocos.

Diante desta situação, temos a seguinte pergunta: existe uma forma
extensível e eficiente para a construção de analisadores léxicos?
Neste capítulo, veremos que a teoria de expressões regulares e
autômatos finitos provêem uma solução eficiente e automatizável para
a criação de analisadores léxicos.

Neste capítulo, vamos revisar a teoria de autômatos, expressões
regulares, a implementação de algoritmos e como estes podem ser
combinados para a criação de analisadores léxicos para uma
especificação de tokens usando expressões regulares.

\section{Autômatos finitos determinísticos}

Nesta seção, vamos apresentar uma implementação simples de autômatos
finitos determinísticos e de operações sobre estes.

\begin{Definition}
  Um autômato finito determinístico M pode ser entendido
  como uma quíntupla \(M = (E,\Sigma, \delta, i, F)\), em que:

  \begin{itemize}
    \item
      \(E\): conjunto finito de estados;
    \item
      \(\Sigma\): alfabeto;
    \item
      \(\delta : E \times \Sigma \to E\), função de transição;
    \item
      \(i \in E\): estado inicial;
    \item
      \(F \subseteq E\): conjunto de estados finais.
  \end{itemize}
\end{Definition}

De maneira simples, o tipo \haskell{DFA} é
parametrizado pelo tipo de estados e possui 3 campos para representar
o estado inicial (campo \haskell{start}), a função de
transição (campo \haskell{delta}) e uma função para
determinar se um estado é final ou não (campo \haskell{finals}).

\begin{minted}{haskell}
data DFA a
  = DFA {
      start :: a
    , delta :: a -> Char -> a
    , finals :: a -> Bool
    }
\end{minted}

Como um exemplo, considere o seguinte autômato que aceita palavras
que terminam em 1.

\begin{figure}[H]
  \begin{tikzpicture}[
      shorten >=1pt,
      node distance=3cm,
      on grid,
      >={Stealth[round]},
      every state/.style={draw, minimum size=12mm, font=\small},
      accepting/.style={double, double distance=1.5pt},
    ]

    \node[state, initial, initial text={}]             (q0) {$q_0$};
    \node[state, accepting, right=of q0] (q1) {$q_1$};

    \path[->]
    (q0) edge [loop above]  node {0} ()
    (q0) edge [bend left]   node [above] {1} (q1)
    (q1) edge [bend left]   node [below] {0} (q0)
    (q1) edge [loop above]  node {1} ();

  \end{tikzpicture}
  \centering
\end{figure}

Esse mesmo autômato pode ser representado pelo seguinte valor em Haskell:

\begin{minted}{haskell}
endsWithOne :: DFA Bool
endsWithOne = DFA False transition id
    where
        transition False '0' = False
        transition False '1' = True
        transition True  '0' = False
        transition True  '1' = True
        transition _ _ = False
\end{minted}

Aqui representamos o conjunto de estados utilizando o tipo
\haskell{Bool}. O estado inicial é codificado por
\haskell{False} e o estado final por
\haskell{True}. A função de transição é representada
por casamento de padrão sobre o estado e o caractere processado. A
função para identificar o estado final é a função identidade,
\haskell{id}.

O cálculo do estado em que a computação de uma string termina em um
AFD é feita percorrendo a string utilizando a função de transição de
forma direta:

\begin{minted}{haskell}
deltaStar :: DFA a -> String -> a
deltaStar m = foldl (delta m) (start m)
\end{minted}

que permite a definição direta de uma função para verificar se uma
palavra pertence ou não a linguagem descrita por um autômato por
testar se o estado retornado por \haskell{deltaStar} é final.

\begin{minted}{haskell}
accept :: DFA a -> String -> Bool
accept m s = finals m (deltaStar m s)
\end{minted}

Evidentemente, em analisadores léxicos devemos utilizar
\textbf{autômatos mínimos}, visto que estes evitam a utilização de
espaço desnecessário por utilizar a menor quantidade possível para
seu conjunto de estados. A implementação de algoritmos de minimização
está disponível no repositório de código associado a este material. O
leitor interessado pode consultar esses detalhes on-line.

\section{Autômatos finitos não determinísticos}

Nesta seção, vamos apresentar uma implementação simples de autômatos
finitos não determinísticos e de operações sobre estes. A seguir,
apresentamos a descrição formal de autômatos não determinísticos.

\begin{Definition}
  Um autômato finito não determinístico M pode ser
  entendido como uma quíntupla \(M = (E,\Sigma, \delta, I, F)\), em que:

  \begin{itemize}
    \item
      \(E\): conjunto finito de estados;
    \item
      \(\Sigma\): alfabeto;
    \item
      \(\delta : E \times \Sigma \to \mathcal{P}(E)\), função de transição;
    \item
      \(I \subseteq E\): conjunto de estados iniciais;
    \item
      \(F \subseteq E\): conjunto de estados finais.
  \end{itemize}
\end{Definition}

Para representação de um autômato não determinístico, vamos
representar de maneira explícita o conjunto de estados do autômato. A
função de transição retorna um conjunto de estados para representar a
possibilidade dos diversos resultados de uma transição em um AFN.

\begin{minted}{haskell}
data NFA a
  = NFA {
      numberOfStates :: Int
    , nfaStart :: Set a
    , nfaDelta :: a -> Char -> Set a
    , nfaFinals :: Set a
    }
\end{minted}

Como um exemplo, considere o seguinte AFN que aceita a linguagem de
palavras que terminam em 00:
\begin{figure}[H]
  \begin{tikzpicture}[
      shorten >=1pt,
      node distance=3cm,
      on grid,
      >={Stealth[round]},
      every state/.style={draw, minimum size=12mm, font=\small},
      accepting/.style={double, double distance=1.5pt},
    ]

    \node[state, initial, initial text={}]                (q0) {$q_0$};
    \node[state, right=of q0]           (q1) {$q_1$};
    \node[state, accepting, right=of q1] (q2) {$q_2$};

    \path[->]
    (q0) edge [loop above] node {0, 1} ()
    (q0) edge              node [above] {0} (q1)
    (q1) edge              node [above] {0} (q2);

  \end{tikzpicture}
  \centering
\end{figure}

este é representado pelo seguinte valor do tipo \haskell{NFA Int}:

\begin{minted}{haskell}
endsWith00 :: NFA Int
endsWith00
    = NFA 3 (Set.singleton 0)
            (\ s c -> case (s,c) of
                        (0, '0') -> Set.fromList [0,1]
                        (0, '1') -> Set.singleton 0
                        (1, '0') -> Set.singleton 2
                        _        -> Set.empty)
            (Set.singleton 2)
\end{minted}

Note que a função de transição é codificada por um par formado por um
estado e o caractere a ser processado pela transição. Os conjuntos de
estados iniciais e finais são representados utilizando a função
\haskell{Set.singleton} que cria um conjunto unitário
com o elemento fornecido como argumento. A definição de aceitação de
uma palavra é similar à de AFDs e está implementada no código de
suporte deste material.

A definição seguinte revisa a equivalência entre AFNs e AFDs.

\begin{Definition}
  Seja $M = (E,\Sigma, \delta, I, F)$ um autômato
  finito não determinístico qualquer. O AFD equivalente a $M$ é
  $M' = (\mathcal{P}(E),\Sigma,\delta',I,F')$, em que:

  \begin{itemize}
    \item
      \(\delta' : \mathcal{P}(E) \times \Sigma \to \mathcal{P}(E)\),
      função de transição;
    \item
      \(I \in \mathcal{P}(E)\): estado inicial;
    \item
      \(F' =\{X \,|\, X \cap F \neq \emptyset\}\): conjunto de estados finais.
  \end{itemize}

  De maneira simples, no AFD equivalente cada estado representa um
  conjunto de estados do AFN. O estado inicial do AFD é apenas o
  conjunto de estados iniciais do AFN e o conjunto de estados finais
  é qualquer conjunto de estados que possui algum estado final do AFN.
\end{Definition}

A partir da definição anterior, podemos codificar uma função que
transforma um AFN qualquer em um AFD equivalente.

\begin{minted}{haskell}
subset :: Ord a => NFA a -> DFA (Set a)
subset m
  = DFA {
      start  = nfaStart m
    , delta  = \ es c ->
        Set.unions (map (flip (nfaDelta m) c) (Set.elems es))
    , finals = \ es -> not (disjoint es (nfaFinals m))
    }
\end{minted}

Observe que a função mostra que se o AFN possui estados de tipo
\haskell{a}, o AFD equivalente terá estados de tipo
\haskell{Set a}, isto é, conjuntos de estados de tipo
\haskell{a}. O estado inicial do AFD é definido como
o estado inicial do AFN, \haskell{nfaStart m}.
Estados finais são definidos por uma função que verifica se o
argumento possui algum estado final do AFN\footnote{A função
  \haskell{disjoint} retorna verdadeiro se a interseção
entre dois conjuntos é vazia.}.

\section{Expressões regulares}

Uma abordagem mais declarativa para análise léxica consiste em
definir os tokens válidos usando expressões regulares e, em seguida,
sintetizar automaticamente o analisador léxico a partir dessas
expressões regulares. Essa abordagem é simples e tem maior
probabilidade de resultar em código eficiente e de fácil manutenção.

Para iniciar, vamos rever a definição da sintaxe de expressões
regulares e de sua semântica.

\begin{Definition}
  Expressões regulares são definidas pela seguinte
  gramática livres de contexto: \[
    \begin{array}{l}
      e \to \emptyset\,|\,\lambda\,|\,a\,|\,e e\,|\,e + e\,|\,e^\star
    \end{array}
  \] Uma expressão regular \(e\) denota um conjunto de strings
  chamado linguagem de \(e\), escrita \(L(e)\), que é definida
  recursivamente como: \[
    \begin{array}{lcl}
      L(\emptyset) & = & \emptyset\\
      L(\lambda) & = & \{\lambda\}\\
      L(a)        & = & \{a\}\\
      L(e_1 e_2)  & = & L(e_1)L(e_2)\\
      L(e_1 + e_2) & = & L(e_1) \cup L(e_2)\\
      L(e^\star)  & = & L(e)^\star
    \end{array}
  \]
\end{Definition}

Existem diversas maneiras para converter uma expressão regular em um
autômato equivalente. Primeiro, vamos revisar uma abordagem que
produz um AFN equivalente por recursão sobre a estrutura da expressão regular.

\begin{Definition}[Construção de Thompson]
  Seja \(e\) uma
  expressão regular arbitrária. O AFN equivalente a \(e\), \(M\), é
  definido por recursão sobre a estrutura de \(e\):

  Caso \(e = \emptyset\): O AFN \(M\) equivalente é

  \begin{figure}[H]
    \begin{tikzpicture}[
        shorten >=1pt,
        node distance=3cm,
        on grid,
        >={Stealth[round]},
        every state/.style={draw, minimum size=12mm, font=\small},
      ]
      \node[state, initial, initial text={}] (q0) {$q_0$};
    \end{tikzpicture}
    \centering
  \end{figure}

  Caso $e = \lambda$: O AFN \(M\) equivalente é
  \begin{figure}[H]
    \begin{tikzpicture}[
        shorten >=1pt,
        node distance=3cm,
        on grid,
        >={Stealth[round]},
        every state/.style={draw, minimum size=12mm, font=\small},
      ]
      \node[state, initial,accepting, initial text={}] (q0) {$q_0$};
    \end{tikzpicture}
    \centering
  \end{figure}

  Caso \(e = a\): O AFN \(M\) equivalente é

  \begin{figure}[H]
    \begin{tikzpicture}[
        shorten >=1pt,
        node distance=3cm,
        on grid,
        >={Stealth[round]},
        every state/.style={draw, minimum size=12mm, font=\small},
        accepting/.style={double, double distance=1.5pt},
      ]

      % Estados
      \node[state, initial, initial text={}]                (q0) {$q_0$};
      \node[state, accepting, right=of q0] (q1) {$q_1$};

      % Transição
      \path[->]
      (q0) edge node [above] {a} (q1);

    \end{tikzpicture}
    \centering
  \end{figure}

  Caso \(e = e_1 + e_2\): O AFN \(M\) equivalente é

  \begin{figure}[H]
    \begin{tikzpicture}[
        shorten >=1pt,
        node distance=2.5cm,
        on grid,
        >={Stealth[round]},
        every state/.style={draw, minimum size=12mm, font=\small},
        accepting/.style={double, double distance=1.5pt},
      ]

      \node[state, initial, initial text={}] (q0) {$q_0$};

      \node[state, above right=2cm and 2.5cm of q0] (q1) {$q_1$};
      \node[state, right=of q1]                      (q2) {$q_2$};

      \node[state, below right=2cm and 2.5cm of q0] (q3) {$q_3$};
      \node[state, right=of q3]                      (q4) {$q_4$};

      \node[state, accepting, below right=2cm and 2.5cm of q2] (qf) {$q_f$};

      \begin{pgfonlayer}{background}
        \node[ellipse, draw=blue!70, fill=blue!10, thick,
          inner sep=8pt, fit=(q1)(q2),
        label={[blue!70, font=\footnotesize]above:$N(e_1)$}] {};
        \node[ellipse, draw=red!70, fill=red!10, thick,
          inner sep=8pt, fit=(q3)(q4),
        label={[red!70, font=\footnotesize]below:$N(e_2)$}] {};
      \end{pgfonlayer}

      \path[->]
      (q0) edge node [above left]  {$\lambda$} (q1)
      (q0) edge node [below left]  {$\lambda$} (q3)
      (q1) edge [blue!70, thick] node [above, text=black] {$e_1$} (q2)
      (q3) edge [red!70, thick]  node [above, text=black] {$e_2$} (q4)
      (q2) edge node [above right] {$\lambda$} (qf)
      (q4) edge node [below right] {$\lambda$} (qf);

    \end{tikzpicture}
    \centering
  \end{figure}

  em que N($e_1$) é o AFN equivalente a $e_1$ e N($e_2$)
  é o AFN equivalente a $e_2$.

  Caso \(e = e_1 e_2\): O AFN \(M\) equivalente é

  \begin{figure}[H]
    \begin{tikzpicture}[
        shorten >=1pt,
        node distance=3cm,
        on grid,
        >={Stealth[round]},
        every state/.style={draw, minimum size=12mm, font=\small},
        accepting/.style={double, double distance=1.5pt},
      ]

      \node[state, initial, initial text={}]    (q0) {$q_0$};
      \node[state, right=of q0] (q1) {$q_1$};

      \node[state, right=of q1] (q2) {$q_2$};
      \node[state, accepting, right=of q2] (q3) {$q_3$};

      \begin{pgfonlayer}{background}
        \node[ellipse, draw=blue!70, fill=blue!10, thick,
          inner sep=10pt, fit=(q0)(q1),
        label={[blue!70, font=\footnotesize]above:$N(e_1)$}] {};
        \node[ellipse, draw=red!70, fill=red!10, thick,
          inner sep=10pt, fit=(q2)(q3),
        label={[red!70, font=\footnotesize]above:$N(e_2)$}] {};
      \end{pgfonlayer}

      \path[->]
      (q0) edge [blue!70, thick]  node [above, text=black] {$e_1$} (q1)
      (q1) edge                   node [above] {$\lambda$}      (q2)
      (q2) edge [red!70, thick]   node [above, text=black] {$e_2$} (q3);
    \end{tikzpicture}
    \centering
  \end{figure}

  em que N($e_1$) é o AFN equivalente a $e_1$ e N($e_2$) é o
  AFN equivalente a $e_2$.

  Caso $e = e_1^\star$: O AFN \(M\) equivalente é

  \begin{figure}[H]
    \begin{tikzpicture}[
        shorten >=1pt,
        node distance=3cm,
        on grid,
        >={Stealth[round]},
        every state/.style={draw, minimum size=12mm, font=\small},
        accepting/.style={double, double distance=1.5pt},
      ]

      \node[state, initial, initial text={}] (q0) {$q_0$};

      \node[state, right=of q0]  (q1) {$q_1$};
      \node[state, right=of q1]  (q2) {$q_2$};

      \node[state, accepting, right=of q2] (qf) {$q_f$};

      \begin{pgfonlayer}{background}
        \node[ellipse, draw=violet!70, fill=violet!10, thick,
          inner sep=10pt, fit=(q1)(q2),
        label={[violet!70, font=\footnotesize]above:$N(e_1)$}] {};
      \end{pgfonlayer}

      \path[->]
      (q0) edge node [above] {$\lambda$} (q1)
      (q1) edge [violet!70, thick] node [above, text=black] {$e_1$} (q2)
      (q2) edge node [above] {$\lambda$} (qf)
      (q2) edge [bend left=50] node [below] {$\lambda$} (q1)
      (q0) edge [bend left=50] node [above] {$\lambda$} (qf);

    \end{tikzpicture}
    \centering
  \end{figure}

  em que N($e_1$) é o AFN equivalente a $e_1$.
\end{Definition}

\section{Gerando AFNs a partir de ERs.}

Nesta seção, vamos apresentar uma implementação Haskell da construção
de Thompson. O primeiro passo é representar expressões regulares
utilizando um tipo de dados Haskell, conforme a seguir:

\begin{minted}{haskell}
data Regex
  = Empty
  | Lambda
  | Chr Char
  | Regex :+: Regex
  | Regex :@: Regex
  | Star Regex
\end{minted}

Para a criação de AFNs a partir da expressão regular, temos que
definir qual tipo representará os estados do AFN. Por questão de
simplicidade, vamos utilizar o tipo \haskell{Int}
para representar estados. Dessa forma, os AFNs produzidos terão tipo
\haskell{NFA Int}.

Iniciamos apresentando o AFN equivalente a expressão regular
\haskell{Empty}:

\begin{minted}{haskell}
emptyNFA :: NFA Int
emptyNFA
  = NFA 0 Set.empty (\ _ _ -> Set.empty) Set.empty
\end{minted}

que consiste de um AFN que não possui estados e nem transições. Dessa
forma, aceita a linguagem vazia\footnote{Optamos por não representar
  como um AFN de um único estado não final para reduzir o número de
estados que seriam considerados inúteis posteriormente.}. O AFN que
aceita a linguagem formada apenas pela palavra vazia é:

\begin{minted}{haskell}
lambdaNFA :: NFA Int
lambdaNFA
  = NFA 1 one (\ _ _ -> Set.empty) one
    where
      one = Set.singleton 1
\end{minted}

que denota um AFN de um único estado inicial e final, sem transições.
A função \haskell{chrNFA} constrói um AFN que aceita
um caractere recebido como argumento.

\begin{minted}{haskell}
chrNFA :: Char -> NFA Int
chrNFA c
  = NFA 2 zero f one
    where
      zero = Set.singleton 0
      one = Set.singleton 1
      err = Set.empty
      f 0 x = if c == x then one
              else err
      f _ _ = err
\end{minted}

Este AFN possui dois estados: 0, que é inicial e não final e 1 que é
final. A função de transição produz como resultado um conjunto vazio,
exceto quando a origem da transição é o estado 0 e o caractere
processado é igual a \haskell{c}, que é o argumento
de \haskell{chrNFA}.

Antes de apresentar funções para produzir AFNs para união,
concatenação e fecho de Kleene, devemos garantir que AFNs não possuam
estados com o mesmo \textbf{nome} (i.e.~mesmo valor de tipo
\haskell{Int}). Para resolver esse problema, vamos
definir uma função que soma um valor \haskell{n} a
cada elemento de um conjunto de inteiros.

\begin{minted}{haskell}
shift :: Int -> Set Int -> Set Int
shift n = Set.fromAscList . map (+ n) . Set.toAscList
\end{minted}

A função \haskell{shift} converte um conjunto de
entrada (tipo \haskell{Set Int}) para uma lista de
inteiros, usando \haskell{Set.toAscList}, e então
soma \haskell{n} a cada elemento da lista produzida
para convertê-la em um conjunto novamente utilizando a função
\haskell{Set.fromAscList}.

A função que calcula o AFN que aceita a união da linguagem de dois
AFNs fornecidos como argumentos é definida como:

\begin{minted}{haskell}
unionNFA :: NFA Int -> NFA Int -> NFA Int
unionNFA m1 m2
  = NFA {
      numberOfStates = n1 + n2
    , nfaStart = Set.union (nfaStart m1) (shift n1 (nfaStart m2))
    , nfaDelta = f
    , nfaFinals = Set.union (nfaFinals m1) (shift n1 (nfaFinals m2))
    }
    where
      n1 = numberOfStates m1
      n2 = numberOfStates m2
      f s c = if s < n1 then nfaDelta m1 s c
              else shift n1 (nfaDelta m2 (s - n1) c)
\end{minted}

A função \haskell{unionNFA} implementa a união de
dois AFNs combinando seus conjuntos de estados de forma disjunta.
Para evitar colisões, ela preserva a estrutura do primeiro autômato
\haskell{m1} e utiliza uma função de deslocamento
\haskell{shift} para renomear os estados do segundo
(\haskell{m2}), garantindo que todos os novos estados
sejam únicos. O resultado é um novo NFA onde os estados iniciais, as
transições e os estados de aceitação são a união das partes
correspondentes de ambos os autômatos originais, permitindo que o AFN
resultante reconheça a linguagem composta pela união das linguagens
aceitas por \haskell{m1} e \haskell{m2}.

A seguir, apresentamos a função que calcula o AFN que aceita
concatenação das linguagem aceitas por dois AFNs fornecidos como argumento.

\begin{minted}{haskell}
concatNFA :: NFA Int -> NFA Int -> NFA Int
concatNFA m1 m2
  = NFA {
      numberOfStates = n1 + n2
    , nfaStart = newStart
    , nfaDelta = newDelta
    , nfaFinals = newFinals
    }
    where
      n1 = numberOfStates m1
      n2 = numberOfStates m2
      start1 = nfaStart m1
      final1 = nfaFinals m1
      final2 = nfaFinals m2
      newStart = if disjoint start1 final1
                 then start1
                 else Set.union start1 (shift n1 final1)
      newFinals = shift n1 final2
      newDelta e c = if e < n1 then
                       if disjoint (nfaDelta m1 e c) final1
                       then nfaDelta m1 e c
                       else Set.union (nfaDelta m1 e c) (shift n1 start2)
                     else shift n1 (nfaDelta m2 (e - n1) c)
\end{minted}

Assim como na construção da união de dois AFNs, utilizamos a função
\haskell{shift} para garantir que estados de cada um
dos AFNs de entrada não possuam o mesmo nome. Para isso, utilizamos o
renomeamento nas definições da função de transição e nos conjuntos de
estados iniciais e finais. A implementação da função de transição
\haskell{newDelta} decide o fluxo da computação com
base no índice do estado atual: se o estado pertence ao primeiro AFN
(\haskell{e < n1}), a função executa a transição
original de \haskell{m1}, mas utiliza a verificação
\haskell{disjoint} para identificar se o destino
atinge um estado de aceitação; em caso positivo, ela funde esse
destino aos estados iniciais para permitir o início imediato do
processamento da segunda linguagem, simulando a conexão entre os
autômatos sem recorrer a transições vazias. Caso o estado já pertença
ao segundo AFN (\haskell{e >= n1}), a função
simplesmente associa a transição original de
\haskell{m2} para o novo intervalo de índices
utilizando \haskell{shift}, garantindo que, uma vez
iniciada a segunda parte da cadeia, o autômato permaneça restrito a
\haskell{m2}.

A função \haskell{starNFA} implementa o fecho de
Kleene ao produzir um AFN, a partir de \haskell{m1},
que aceita zero ou mais repetições da linguagem original, sem
adicionar novos estados. Para permitir a aceitação da string vazia e
o reinício do ciclo, a função redefine os estados de aceitação
(\haskell{nfaFinals}) como sendo os próprios estados
iniciais (\haskell{nfaStart m1}) e modifica a função
de transição (\haskell{newDelta}): sempre que uma
transição original de \haskell{m1} atingiria um
estado final, ela é inclui os estados iniciais usando
\haskell{Set.union}. Esse detalhe cria um laço que
permite ao autômato processar múltiplas palavras consecutivas,
tratando o alcance de um estado final como um gatilho para reiniciar
a computação a partir de algum dos estados iniciais.

\begin{minted}{haskell}
starNFA :: NFA Int -> NFA Int
starNFA m1
  = NFA {
      numberOfStates = numberOfStates m1
    , nfaStart = nfaStart m1
    , nfaDelta = newDelta
    , nfaFinals = nfaStart m1
    }
    where
      newDelta e c
        = let r = nfaDelta m1 e c
          in if disjoint r (nfaFinals m1)
             then r
             else Set.union r (nfaStart m1)
\end{minted}

De posse de todas essas funções auxiliares, a conversão de uma
expressão regular em um AFN equivalente é definida por recursão sobre
a estrutura da expressão:

\begin{minted}{haskell}
thompson :: Regex -> NFA Int
thompson Empty = emptyNFA
thompson Lambda = lambdaNFA
thompson (Chr c) = chrNFA c
thompson (e1 :+: e2)
  = unionNFA (thompson e1) (thompson e2)
thompson (e1 :@: e2)
  = concatNFA (thompson e1) (thompson e2)
thompson (Star e1)
  = starNFA (thompson e1)
\end{minted}

Especificações léxicas de uma linguagem são feitas utilizando
diversas expressões regulares. A função a seguir, combina as
implementações descritas até o momento para produzir o AFD
equivalente a uma lista de expressões regulares:

\begin{minted}{haskell}
lexer :: [Regex] -> DFA (Set Int)
lexer = subset . foldr unionNFA emptyNFA . map thompson
\end{minted}

Primeiro, obtemos um AFN equivalente para cada expressão usando a
função \haskell{thompson}. Em seguida, realizados a
união de todos os AFNs produzidos utilizando a função
\haskell{unionNFA}. Finalmente, obtemos o AFD
equivalente chamando a função \haskell{subset} sobre
o resultado dos passos anteriores.

\section{Critério do maior prefixo}

Autômatos fornecem um maneira eficiente e fundamentada para a criação
de analisadores léxicos. Porém, a mera aceitação de um prefixo da
entrada é suficiente para que um analisador produza um token? Para
entender considere o seguinte trecho de código:

\begin{minted}{java}
if (x > 0) ift = x + 1;
\end{minted}

Como diferenciar o identificador \haskell{ift} da
palavra reservada \haskell{if}? A ideia é que
analisadores léxicos sempre busquem aceitar o \textbf{maior prefixo}
possível. Esse critério é necessário para evitar ambiguidades
semânticas e garantir que a análise reflita a intenção real da
linguagem. Sem esse critério, o analisador léxico poderia fragmentar
precocemente sequências de caracteres que possuem um significado
específico quando combinadas, gerando múltiplos tokens genéricos em
vez de um único token. Com isso, o compilador assegura que
identificadores, palavras-reservadas e operadores compostos sejam
interpretados de forma atômica e correta dentro do contexto gramatical.

Mas isso nos leva a seguinte pergunta: Como reconhecer o maior
prefixo utilizando um AFD? A função a seguir, permite que um AFD de
argumento reconheça o maior prefixo de uma string de entrada.

\begin{minted}{haskell}
longest :: DFA a -> String -> Maybe String
longest m = combine . foldl step (Just "", Nothing, start m)
  where
    step (Just pre, Nothing, e) c
      | finals m (delta m e c)
          = ( Just (c : pre)
            , Just (c : pre)
            , delta m e c)
      | otherwise
          = ( Just (c : pre)
            , Nothing
            , delta m e c)
    step (Just pre, Just pre', e) c
      | finals m (delta m e c)
          = ( Just (c : pre)
            , Just (c : pre')
            , delta m e c)
      | otherwise = ( Nothing
                    , Just pre'
                    , delta m e c)
    step (Nothing, val, e) c = ( Nothing
                               , val
                               , delta m e c)

    combine (_, val, _) = reverse <$> val
\end{minted}

A função \haskell{longest} busca identificar o maior
prefixo de uma string aceita por AFD, percorrendo a entrada enquanto
mantém um acumulador representado como uma tripla formada por:

\begin{itemize}
  \item
    Prefixo atual sendo processado
  \item
    Último prefixo válido (aceito pelo DFA) encontrado até o momento
  \item
    Estado atual do autômato.
\end{itemize}

A lógica de atualização do acumulador é feita pela função
\haskell{step} que lida com três situações: ela
acumula caracteres no prefixo atual enquanto for possível estendê-lo,
salva essa string como o ``melhor resultado'' sempre que atinge um
estado final (\haskell{finals}) e, caso o caminho se
torne inválido para futuras aceitações (indicado pelo valor
\haskell{Nothing} no primeiro elemento da tupla),
interrompe a atualização do prefixo em progresso mas preserva a maior
correspondência já registrada. Ao final, a função
\haskell{combine} recupera esse maior prefixo
armazenado e inverte a ordem da string (corrigindo a inversão causada
pela construção da lista) para retornar a string resultante.

\section{Derivadas de expressões regulares}

A geração de analisadores léxicos baseados na construção de Thompson
é útil para o reconhecimento de expressões regulares em um
compilador, visto que a especificação do token não se altera. No
entanto, expressões regulares são frequentemente usadas em outros
contextos em que a pré-compilação em um AFD não compensa o custo.
Infelizmente, bibliotecas comuns de expressões regulares, como as
utilizadas por Java, Perl e Javascript, apresentam desempenho
exponencial no pior caso, pois empregam backtracking para lidar com o
operador de união.

Uma técnica alternativa consiste em interpretar diretamente a
expressão regular à medida que a entrada é analisada, usando
derivadas de expressões regulares, uma técnica elegante desenvolvida
por Janus Brzozowski. A ideia é que, para qualquer expressão regular
e símbolo de entrada, podemos calcular a expressão regular que aceita
o restante da entrada após o símbolo de entrada ter sido consumido.
Usamos a notação \(\partial(e,a)\) para representar a derivada da
expressão regular \(e\) com respeito ao símbolo de entrada \(a\).

Antes de apresentar a definição de derivadas propriamente ditas,
precisamos de um conceito auxilar: a noção de anulabilidade. Dizemos
que uma expressão regular é \textbf{anulável} se \(\lambda \in
L(e)\). Esse conceito é definido a seguir.

\begin{Definition}
  Dizemos que uma expressão regular é \textbf{anulável}
  se \(\nu(e) = \top\), em que \(\nu(e)\) é definido por recursão
  sobre a estrutura da expressão da seguinte maneira:

  \[
    \begin{array}{lcl}
      \nu(\emptyset) & = & \bot\\
      \nu(\lambda)   & = & \top\\
      \nu(a)          & = & \bot\\
      \nu(e_1 e_2)    & = & \nu(e_1) \land \nu(e_2)\\
      \nu(e_1 + e_2)  & = & \nu(e_1) \lor \nu(e_2)\\
      \nu(e_1^\star) & = & \top\\
    \end{array}
  \]
\end{Definition}

As três primeiras equações mostram que o conjunto vazio não é
anulável, que \(\lambda\) é anulável e que um símbolo não é
considerado anulável. Uma concatenação é anulável somente se ambas as
suas sub-expressões são anuláveis e a união é considerada anulável se
qualquer uma de suas sub-expressões o for anulável. Finalmente, uma
expressão envolvendo o fecho de Kleene é sempre anulável.

Utilizando o conceito de anulabilidade, podemos definir a derivada de
expressões regulares.

\begin{Definition}[Derivada de uma expressão regular]
  Seja \(e\) uma expressão regular qualquer e \(a\in
  \Sigma\). A derivada de \(e\) com respeito a \(a\) é definida como:

  \[
    \begin{array}{lcl}
      \partial(\emptyset,a) & = & \emptyset\\
      \partial(\lambda, a)  & = & \emptyset\\
      \partial(b, a)         & = & \left\{
        \begin{array}{ll}
          \lambda & \textrm{se }a = b \\
          \emptyset & \textrm{caso contrário}
        \end{array}
        \right . \\
        \partial(e_1 + e_2, a) & = & \partial(e_1,a) + \partial(e_2,a)\\
        \partial(e_1 e_2, a)   & = & \left\{
          \begin{array}{ll}
            \partial(e_1, a) e_2 + \partial(e_2, a) & \textrm{se
            }\nu(e_1) = \top\\
            \partial(e_1,a) e_2 & \textrm{caso contrário}
          \end{array}
          \right .\\
          \partial(e_1^\star)   & = & \partial(e_1,a)e_1^\star
        \end{array}
      \]
    \end{Definition}

    A definição da derivada de uma expressão regular em relação a um
    símbolo \$\$ é apresentada de forma recursiva. Para os casos
    base: a derivada da expressão vazia \(\emptyset\) ou da palavra
    vazia \(\lambda\) em relação a qualquer símbolo resulta em
    \(\emptyset\), pois não é possível consumir um caractere a partir
    delas. No caso de um símbolo \(b\), a derivada retorna
    \(\lambda\) se o símbolo consumido \(a\) for igual a \(b\)
    (indicando que o caractere esperado foi encontrado e a expressão
    se reduz à palavra vazia), ou \(\emptyset\) caso contrário. Para
    a união \(e_1 + e_2\), a derivada é simplesmente a união das
    derivadas de cada subexpressão. No caso da concatenação \(e_1
    e_2\), a definição se desdobra: se \(e_1\)\hspace{0pt} for
    anulável (ou seja, pode reconhecer a palavra vazia), a derivada é
    a união entre a concatenação da derivada de \(e_1\) com \(e_2\) e
    a derivada de \(e_2\)\hspace{0pt}; caso contrário, é apenas a
    derivada de \(e_1\)\hspace{0pt} concatenada com \(e_2\). Por fim,
    para o fecho de Kleene \(e^\star\), a derivada é a concatenação
    da derivada de \(e\)\hspace{0pt} com a própria estrela
    \(e^\star\), permitindo que a expressão continue reconhecendo
    múltiplas repetições.

    Para obter um autômato finito determinístico (AFD) a partir de
    uma expressão regular utilizando derivadas, podemos construir os
    estados do autômato como as próprias expressões regulares que
    surgem durante o processo de derivação. A ideia fundamental é
    que, partindo da expressão regular original \(e_0\)\hspace{0pt},
    podemos calcular sua derivada em relação a cada símbolo do
    alfabeto, obtendo novas expressões que representam o estado após
    consumir aquele símbolo. Este processo é repetido para cada nova
    expressão gerada, até que nenhuma nova expressão seja produzida,
    garantindo que o conjunto de estados seja finito.

    \begin{Definition}
      Seja \(e\) uma expressão regular qualquer. O AFD
      equivalente a \(e\) é \(M = (E, \Sigma, \delta, i, F)\), em
      que:
      \begin{itemize}
        \item \(E\): é o conjunto de todas as derivadas para a
          expressão \(e\).
        \item \(\Sigma\): alfabeto de entrada
        \item \(\delta\): função de transição, representada pela própria
          derivada: \(\delta(e,a) = \partial(e,a)\).
        \item \(F = \{e_1 \,|\,
          \nu(e_1) = \top\}\), isto é, consideramos como estados finais
          as expressões regulares anuláveis.
      \end{itemize}
    \end{Definition}

    Um ponto crucial nesta construção é o fato de que o conjunto de
    derivadas de uma expressão é \textbf{finito}. Porém, esse fato é
    verdade apenas se considerarmos algumas simplificações,
    apresentadas a seguir:
    \[
      \begin{array}{ccc}
        e_1 + \emptyset \equiv e_1 &
        \emptyset + e_2 \equiv e_2 &
        e_1\emptyset \equiv \emptyset \\
        \emptyset e_2 \equiv \emptyset &
        e_1 \lambda\equiv e_1 &
        \lambda e_2 \equiv e_2 \\
        \lambda^\star \equiv \lambda &
        \emptyset^\star \equiv \lambda &
        (e^\star)^\star \equiv e^\star
      \end{array}
    \]

    Considerando essas simplificações, é demonstrado que sempre
    obteremos um conjunto finito de derivadas a partir de uma
    expressão regular qualquer.

    \section{Implementação em Haskell}

    De posse das definições, podemos implementar a obtenção de um AFD
    a partir de uma expressão regular de forma direta. Primeiro,
    vamos definir uma função para determinar se uma expressão é ou não anulável.

\begin{minted}{haskell}
nullable :: Regex -> Bool
nullable Empty = False
nullable Lambda = True
nullable (Chr _) = False
nullable (e1 :+: e2) = nullable e1 || nullable e2
nullable (e1 :@: e2) = nullable e1 && nullable e2
nullable (Star _) = True
\end{minted}

    Observe que a implementação é uma tradução direta da notação
    matemática para código Haskell. Em seguida, apresentamos funções
    para realizar as simplificações da união, concatenação e fecho de
    Kleene. Cada uma destas funções realiza as simplificações
    apresentadas anteriormente.

\begin{minted}{haskell}
(.+.) :: Regex -> Regex -> Regex
Empty .+. e' = e'
e .+. Empty  = e
e .+. e'     = e :+: e'

(.@.) :: Regex -> Regex -> Regex
Empty .@. _ = Empty
_ .@. Empty = Empty
Lambda .@. e' = e'
e .@. Lambda = e
e .@. e' = e :@: e'

star :: Regex -> Regex
star Empty = Lambda
star (Star e) = e
star e = Star e
\end{minted}

    De posse destas simplificações e do teste de anulabilidade, a
    implementação da operação de derivada é apenas a tradução da
    notação matemática para código Haskell.

\begin{minted}{haskell}
deriv :: Char -> Regex -> Regex
deriv _ Empty  = Empty
deriv _ Lambda = Empty
deriv a (Chr b)
  | a == b = Lambda
  | otherwise = Empty
deriv a (e1 :+: e2)
  = deriv a e1 .+. deriv a e2
deriv a (e1 :@: e2)
  | nullable e1 = deriv a e1 .@. e2 .+. deriv a e2
  | otherwise   = deriv a e1 .@. e2
deriv a (Star e1)
  = deriv a e1 .@. (star e1)
\end{minted}

    Para construção do AFD utilizando derivadas, precisamos obter o
    conjunto de todas as derivadas de uma expressão regular. A função
    \haskell{enumerateDerivs} implementa um algoritmo
    para enumerar todas as expressões regulares alcançáveis a partir
    de uma expressão inicial (\haskell{initial})
    através da aplicação sucessiva da operação de derivada em relação
    aos símbolos do alfabeto.

\begin{minted}{haskell}
enumerateDerivs :: Regex -> [Regex]
enumerateDerivs initial = go [initial] [initial]
  where
    go [] acc = acc
    go (state:rest) acc =
        let chars = symbols initial
            newStates = nub [deriv c state | c <- chars]
            unseen = filter (`notElem` acc) newStates
        in go (rest ++ unseen) (acc ++ unseen)
\end{minted}

    A função \haskell{enumerateDerivs} realiza uma
    busca sistemática em largura sobre o espaço de estados
    (expressões regulares), utilizando duas listas: a primeira
    (state:rest) funciona como uma fila de estados a serem
    processados, enquanto a segunda (acc) acumula todos os estados já
    visitados. Para cada estado atual, a função obtém todos os
    símbolos do alfabeto a partir da expressão inicial (usando
    \haskell{symbols initial}), calcula a derivada
    para cada um deles com \haskell{deriv c state}, e
    remove duplicidades com \haskell{nub}. As novas
    expressões geradas que ainda não estão no acumulador

    (\haskell{filter (notElem acc) newStates})

    são então adicionadas tanto ao final da fila de processamento quanto
    ao acumulador, garantindo que todos os estados alcançáveis sejam
    explorados. O processo continua até que a fila se esgote,
    retornando a lista completa de expressões regulares que formam os
    estados do autômato.

    Finalmente, a construção do autômato a partir da expressão
    regular é feita pela função \haskell{derivDFA}:

\begin{minted}{haskell}
derivDFA :: Regex -> DFA Regex
derivDFA e
  = DFA {
      start = initial
    , delta = delta'
    , finals = nullable
    }
  where
    initial = simp e
    states = enumerateDerivs initial

    delta' state char =
        let e' = deriv char state
        in if e' `elem` states then e' else Empty
\end{minted}

    O estado inicial é dado pela simplificação da expressão regular
    de entrada e os estados finais são exatamente as expressões
    regulares anuláveis. Finalmente, a função de transição é a
    derivada com respeito ao caractere considerado.

    \section{Conclusão}

    Neste capítulo, apresentamos os fundamentos teóricos e
    implementações práticas para a construção automatizada de
    analisadores léxicos. Exploramos autômatos finitos
    determinísticos e não determinísticos, a construção de Thompson
    para conversão de expressões regulares em AFNs, e o algoritmo
    para converter AFNs em AFDs. Discutimos o critério do maior
    prefixo como solução para ambiguidades na identificação de
    tokens. Por fim, introduzimos a abordagem de derivadas de
    expressões regulares, que permite construir AFDs diretamente a
    partir de expressões regulares de forma elegante e eficiente.
